\documentclass[11pt,a4paper]{article}

\usepackage[polish]{babel}
\usepackage[utf8]{inputenc}
\usepackage{polski}
\usepackage[T1]{fontenc}
\usepackage{indentfirst}
\usepackage{wrapfig}    % for wrapping figures, tables
\usepackage{isotope}
\usepackage{hyperref}
\usepackage{amsmath}
\usepackage{amssymb}
\usepackage{bm}
\usepackage{gensymb}
\usepackage{epsfig}
\usepackage{graphics}
\usepackage[shortlabels]{enumitem}
\usepackage{xspace}
\xspaceaddexceptions{[]\{\}}

%\def\hidesolutions{}  %← show solutions. Comment out to show solutions

%
%
%fixpagesize
\pagestyle{empty}
\addtolength{\textwidth}{4cm}
\addtolength{\textheight}{6cm}
\addtolength{\evensidemargin}{-3cm}
\addtolength{\oddsidemargin}{-2cm}
\addtolength{\topmargin}{-3cm}
\parindent=0cm

%
%
% Changes figure placing algorithm
\renewcommand{\topfraction}{1}       % maximal fraction of a page allowed for figures
\renewcommand{\textfraction}{0.15}   % minimal number of text for figure-text shared pages
\renewcommand{\floatpagefraction}{0.95} % if two above does not help, this could do the job 
                                        % must be: floatpagefraction < topfraction !!!!
\renewcommand{\textfraction}{0} % minimum fraction of page, which must be  devoted to text
\renewcommand{\topfraction}{1}  % maximum fraction at top, which can be occupied whit floats
\setcounter{totalnumber}{400}   % increase the number of floats for one page
\setcounter{topnumber}{200}     % at all/top/bottom.
\setcounter{bottomnumber}{200}  %

%
%
%small distance in list/item/enum for enumitem package
\setlist[itemize,enumerate]{topsep=0em}
\setlist{noitemsep}

%Nuclear notations: usage \nucl{235}{92}{U}. Math mode optional
\newcommand{\nucl}[3]{\ensuremath{
  \phantom{\ensuremath{^{#1}_{#2}}}
  \llap{\ensuremath{^{#1}}}
  \llap{\ensuremath{_{\rule{0pt}{.75em}#2}}}
  \mbox{#3} } 
} 

\newcounter{zadanie}\newcommand{\zadanie}[1][]{\addtocounter{zadanie}{1} ~\\  {\bf \emph{Exercise \arabic{zadanie} #1 }} \\}
\newcounter{zaddom}\newcommand{\zaddom}[1][]{\addtocounter{zaddom}{1} ~\\  {\bf \emph{Homework \arabic{zaddom} #1 }} \\}

%dynamic solutions hiding
\usepackage{comment}

\newif\ifhidesolutions
\ifdefined\hidesolutions
  \hidesolutionstrue
\else
  \hidesolutionsfalse
\fi

\ifhidesolutions
  \excludecomment{solution}
\else
  \newenvironment{solution}{\vspace{1cm} \par\textbf{Solution.} \vspace{1cm}}{}
\fi
%%%%%%%%%%%%%%%%%%%%%%%%%%%%%%%%%%%%%%%%%%
\begin{document}         

%%%%%%%%%%%%%%%%%%%%%%%%%%%%%%%%%%%%%%%%%%%%%%%%%%%%%
\begin{centering}
  {\bf {\Large
 Introduction to Experimental Particle Physics I \\
  Set III: Lorentz transformation, cross section, decay width, branching ratio
  }} \\  
\end{centering}

\vspace{1cm}
%%%%%%%%%%%%%%%%%%%%%%%%%%%%%%%%%%%%%%%%%%%%%%%%%%%%%%%
\zadanie[Boosts and rotations]

Calculate numerical values of Lorentz transformation matrix composed of two boosts
along $\vec{\beta}_1 = (0.6,0,0)$ and $\vec{\beta}_2 = (0, 0.6,0)$.
Calculate the same for a single, combined, boost along $\vec{\beta}_3 = (0.6, 0.48, 0)$.
\begin{itemize}
  \item make numerical checks if resulting matrices represent a Lorentz transformation, 
  \item explain why combined boost is along $\vec{\beta}_3 = (0.6, 0.48, 0)$ and not 
  along $\vec{\beta}_1 + \vec{\beta}_2 = (0.6, 0.6, 0)$,
  \item compare the results and discuss the difference.
\end{itemize}

\vspace{0.2cm}
Use Numpy to perform matrix multiplications, and submit solution in form of a short ipynb 
notebook with calculations and checks.

\begin{solution}

Lorentz matrix for a boost along $\vec{\beta}_1 = (0.6,0,0)$ is:
\begin{align*}
  \Lambda_1 = \begin{pmatrix}
    \gamma_1 & -\beta_1 \gamma_1 & 0 & 0 \\
    -\beta_1 \gamma_1 & \gamma_1 & 0 & 0 \\
    0 & 0 & 1 & 0 \\
    0 & 0 & 0 & 1
  \end{pmatrix} = \begin{pmatrix}
    1.25 & -0.75 & 0 & 0 \\
    -0.75 & 1.25 & 0 & 0 \\
    0 & 0 & 1 & 0 \\
    0 & 0 & 0 & 1
  \end{pmatrix}
\end{align*}

Matrix validity test: $\Lambda_1^T g \Lambda_1 = g$:
\begin{align*}
  \Lambda_1^T g \Lambda_1 = \begin{pmatrix}
    1.25 & -0.75 & 0 & 0 \\
    -0.75 & 1.25 & 0 & 0 \\
    0 & 0 & 1 & 0 \\
    0 & 0 & 0 & 1
  \end{pmatrix} \begin{pmatrix}
    1 & 0 & 0 & 0 \\
    0 & -1 & 0 & 0 \\
    0 & 0 & -1 & 0 \\
    0 & 0 & 0 & -1
  \end{pmatrix} \begin{pmatrix}
    1.25 & -0.75 & 0 & 0 \\
    -0.75 & 1.25 & 0 & 0 \\
    0 & 0 & 1 & 0 \\
    0 & 0 & 0 & 1
  \end{pmatrix} = \begin{pmatrix}
    1 & 0 & 0 & 0 \\
    0 & -1 & 0 & 0 \\
    0 & 0 & -1 & 0 \\
    0 & 0 & 0 & -1
  \end{pmatrix} = g
\end{align*}

Determinant of $\Lambda_1$ is $det(\Lambda_1) = 1.25 \cdot 1.25 - (-0.75) \cdot (-0.75) = 1$.


Lorentz matrix for a boost along $\vec{\beta}_2 = (0,0.6,0)$ is:
\begin{align*}
  \Lambda_2 = \begin{pmatrix}
    \gamma_2 & 0 & -\beta_2 \gamma_2 & 0 \\
    0 & 1 & 0 & 0 \\
    -\beta_2 \gamma_2 & 0 & \gamma_2 & 0 \\
    0 & 0 & 0 & 1
  \end{pmatrix} = \begin{pmatrix}
    1.25 & 0 & -0.75 & 0 \\
    0 & 1 & 0 & 0 \\
    -0.75 & 0 & 1.25 & 0 \\
    0 & 0 & 0 & 1
  \end{pmatrix}
\end{align*}

\vspace{0.5cm}

Metric conservation test: $\Lambda_2^T g \Lambda_2 = g$:
\begin{align*}
  \Lambda_2^T g \Lambda_2 = \begin{pmatrix}
    1.25 & 0 & -0.75 & 0 \\
    0 & 1 & 0 & 0 \\
    -0.75 & 0 & 1.25 & 0 \\
    0 & 0 & 0 & 1
  \end{pmatrix} \begin{pmatrix}
    1 & 0 & 0 & 0 \\
    0 & -1 & 0 & 0 \\
    0 & 0 & -1 & 0 \\
    0 & 0 & 0 & -1
  \end{pmatrix} \begin{pmatrix}
    1.25 & 0 & -0.75 & 0 \\
    0 & 1 & 0 & 0 \\
    -0.75 & 0 & 1.25 & 0 \\
    0 & 0 & 0 & 1
  \end{pmatrix} = \begin{pmatrix}
    1 & 0 & 0 & 0 \\
    0 & -1 & 0 & 0 \\
    0 & 0 & -1 & 0 \\
    0 & 0 & 0 & -1
  \end{pmatrix} = g
\end{align*}

Determinant test: $det(\Lambda_2) = 1.25 \cdot 1.25 - (-0.75) \cdot (-0.75) = 1$.

\vspace{0.5cm}

Matrix for consecutive boosts along $\vec{\beta}_1$ and $\vec{\beta}_2$ is:
\begin{align*}
  \Lambda_{12} = \Lambda_2 \Lambda_1 = \begin{pmatrix}
    1.5625 & -0.9375 & -0.9375 & 0 \\
    -0.75 & 1.25 & 0 & 0 \\
    -0.9375 & 0.5625 & 1.5625 & 0 \\
    0 & 0 & 0 & 1
  \end{pmatrix}
\end{align*}

Matrix for a single boost along $\vec{\beta}_3 = (0.6, 0.48, 0)$ is:
\begin{align*}
  \Lambda_3 = \begin{pmatrix}
    \gamma_3 & -\beta_{3x} \gamma_3 & -\beta_{3y} \gamma_3 & 0 \\
    -\beta_{3x} \gamma_3 & 1 + (\gamma_3 - 1) \frac{\beta_{3x}^2}{\beta_3^2} & (\gamma_3 - 1) \frac{\beta_{3x}\beta_{3y}}{\beta_3^2} & 0 \\
    -\beta_{3y} \gamma_3 & (\gamma_3 - 1) \frac{\beta_{3x}\beta_{3y}}{\beta_3^2} & 1 + (\gamma_3 - 1) \frac{\beta_{3y}^2}{\beta_3^2} & 0 \\
    0 & 0 & 0 & 1
  \end{pmatrix} = \begin{pmatrix}
    1.5625 & -0.9375 & -0.75 & 0 \\
    -0.9375 & 1.343 & 0.2744 & 0 \\
    -0.75 & 0.2744 & 1.2195 & 0 \\
    0 & 0 & 0 & 1
  \end{pmatrix}
\end{align*}

The combined boost along $\vec{\beta}_3$ takes into account the relativistic velocity addition formula, 
which results in a different effective velocity than the naive sum of $\vec{\beta}_1$ and $\vec{\beta}_2$.
The boost along combined direction effects in a pure boost and a rotation (Thomas-Wigner rotation) in the plane of 
the two boosts, which is why the resulting matrix $\Lambda_3$ differs from $\Lambda_{12}$.

\end{solution}
%%%%%%%%%%%%%%%%%%%%%%%%%%%%%%%%%%%%%%%%%%%%%%%%%%%%%%%
\ifdefined\hidesolutions
  \newpage
\fi
%%%%%%%%%%%%%%%%%%%%%%%%%%%%%%%%%%%%%%%%%%%%%%%%%%%%%%%
\zadanie[Two body decay phase space]

Calculate the width for a two body decay of a particle with mass $M$ into two 
particles with masses $m_1$ and $m_2$. Assume general form of decay matrix element 
$|\mathcal{M}|^2$, and leave this as it is in the final formula. Consider two cased:
\begin{itemize}
  \item final state particles are different, $m_1 \neq m_2$,
  \item final state particles are identical, $m_1 = m_2$.
\end{itemize}

\begin{solution}
  The width of the decay is given by:
  \begin{align*}
    \Gamma = \frac{1}{2M} \int \frac{d^3p_1}{(2\pi)^3 2E_1} \frac{d^3p_2}{(2\pi)^3 2E_2} (2\pi)^4 \delta^4(p - p_1 - p_2) |\mathcal{M}|^2
  \end{align*}
  
  where $p$ is the four-momentum of the decaying particle, and $p_1$ and $p_2$ are the four-momenta of the decay products.
  We consider the rest frame of the decaying particle, where $p = (M, 0, 0, 0)$.

  \begin{align*}
    \Gamma = & \frac{1}{2M} \int \frac{d^3p_1}{(2\pi)^3 2E_1} \frac{d^3p_2}{(2\pi)^3 2E_2} (2\pi)^3 \delta^4(\vec{p} - \vec{p_1} - \vec{p_2}) 
    (2\pi) \delta(M - E_1 - E_2) |\mathcal{M}|^2 = \\
    & | \text{integrate over } d^3p_2 \text{ using the delta function } \delta^3(\vec{p} - \vec{p_1} - \vec{p_2}) | = \\
    & \frac{1}{2M} \int \frac{d^3p_1}{(2\pi)^3 2E_1} \frac{1}{(2\pi)^3 2E_2} (2\pi)^4 \delta(M - E_1 - E_2) |\mathcal{M}|^2 = \\
    & \frac{1}{2M} \int \frac{p_{1}^{2}dp_1 d\Omega_1}{(2\pi)^3 2E_1} \frac{1}{(2\pi)^3 2E_2} (2\pi)^4 \delta(M - E_1 - E_2) |\mathcal{M}|^2 = \\
    & | \delta(M - E_1 - E_2) = \delta(M - \sqrt{|\vec{p}_1|^2 + m_1^2} + \sqrt{|\vec{p}_1|^2 + m_2^2}) =  \\
    & \delta(p_{1} - p^{*}) \left( \frac{p_{1}}{E_{1}} + \frac{p_{1}}{E_{2}} \right)^{-1} = 
      \delta(p_{1} - p^{*}) \frac{E_{1} E_{2}}{p_{1} (E_{1} + E_{2})} = \\
    & || E_1 + E_2 = M || = 
     \delta(p_{1} - p^{*}) \frac{E_{1} E_{2}}{p_{1} M} =  \\
    & \frac{1}{2M} \frac{(2\pi)^4 p^{*}}{(2\pi)^6 4M} \int  |\mathcal{M}|^2 d\Omega = \\
    & | \text{assume particle is decaying isotropically, then} |\mathcal{M}| \\
    & \text{does not depend on the direction of the decay products} | = \\
    &\frac{1}{2M} \frac{(2\pi)^4 p^{*}}{(2\pi)^6 4M} |\mathcal{M}|^2 4\pi = 
     \frac{p^{*}}{8\pi M^2} |\mathcal{M}|^2
  \end{align*}

For identical particles in the final state, we need to include a factor of $\frac{1}{2}$ 
to account for the indistinguishability of the particles, which means we should integrate only 
over half of the solid angle. Therefore, the width for identical particles is:
\begin{align*}
  \Gamma_{identical} = \frac{1}{2} \frac{p^{*}}{8\pi M^2} |\mathcal{M}|^2 = \frac{p^{*}}{16\pi M^2} |\mathcal{M}|^2
\end{align*}

\end{solution}
%%%%%%%%%%%%%%%%%%%%%%%%%%%%%%%%%%%%%%%%%%%%%%%%%%%%%%%
%%%%%%%%%%%%%%%%%%%%%%%%%%%%%%%%%%%%%%%%%%%%%%%%%%%%%%%
\end{document}
